% Bagian: computability
% Bab: recursive-functions
% Subbagian: general-recursive-functions

\documentclass[../../../include/open-logic-section]{subfiles}

\begin{document}

\olfileid[id]{cmp}{rec}{gen}
\olsection{Fungsi Rekursif Umum}

Ada cara lain untuk memperoleh suatu himpunan fungsi total. Suatu fungsi total
$f(x,\vec z)$ disebut \emph{reguler} jika, bagi setiap barisan bilangan asli
$\vec z$, terdapat suatu $x$ sedemikian sehingga $f(x,\vec z)=0$. Dengan kata
lain, fungsi reguler tepat merupakan fungsi yang dapat dikenai pencarian tak
berbatas dan menghasilkan fungsi total. Secara konservatif, kita dapat
membatasi pencarian tak berbatas pada fungsi reguler:

\begin{defn}
\ollabel{defn:general-recursive}
Himpunan \emph{fungsi rekursif umum} ialah himpunan terkecil fungsi pada
bilangan asli (dengan berbagai aritas) yang memuat nol, penerus, dan proyeksi,
serta tertutup terhadap komposisi, rekursi primitif, dan pencarian tak berbatas
yang diterapkan pada fungsi \emph{reguler}.
\end{defn}

Jelas bahwa setiap fungsi rekursif umum bersifat total. Perbedaan antara
\olref{defn:general-recursive} dan \olref[par]{defn:recursive-fn} ialah bahwa
dalam definisi kedua, kita diperbolehkan menggunakan fungsi rekursif parsial
selama proses; satu-satunya syarat ialah bahwa fungsi yang akhirnya diperoleh
bersifat total. Jadi, kata ``umum'', suatu peninggalan historis, sebenarnya
kurang tepat; sepintas, \olref{defn:general-recursive} justru \emph{kurang}
umum daripada \olref[par]{defn:recursive-fn}. Namun, untungnya, perbedaan itu
hanya semu; meskipun definisinya berbeda, himpunan fungsi rekursif umum dan
himpunan fungsi rekursif sama persis.

\end{document}
